% 06_tablas/matriz_normativa.tex
% Tabla 2. Matriz de instrumentos normativos aplicables al sector minero y al derecho ambiental
% Capítulo II, sección 2.5 — Marco legal
% Usa longtable para permitir extensión en varias páginas

\begingroup
\small
\begin{longtable}{|p{0.03\textwidth}|p{0.22\textwidth}|p{0.18\textwidth}|p{0.22\textwidth}|p{0.24\textwidth}|}

\caption{Matriz de instrumentos normativos aplicables al sector minero y al derecho ambiental ecuatoriano}
\label{tab:normativa} \\

\hline
\textbf{N.°} & \textbf{Instrumento normativo} & \textbf{Tipo / Registro Oficial} & \textbf{Artículos clave para la tesis} & \textbf{Relevancia analítica} \\
\hline
\endfirsthead

\multicolumn{5}{l}{\small\textit{(continuación de la Tabla \ref*{tab:normativa})}} \\[2pt]
\hline
\textbf{N.°} & \textbf{Instrumento normativo} & \textbf{Tipo / Registro Oficial} & \textbf{Artículos clave para la tesis} & \textbf{Relevancia analítica} \\
\hline
\endhead

\hline
\endfoot

1 &
Constitución de la República del Ecuador (2008) &
Norma suprema. R.O.\ No.\ 449, 20-oct-2008 &
Arts.\ 14, 57.7, 71--74, 82, 316, 395--399, 407--409 &
Parámetro de validez constitucional de toda legislación minera y ambiental. Reconoce a la naturaleza como sujeto de derechos y establece principios ambientales vinculantes. \\
\hline

2 &
Código Orgánico del Ambiente (COA) (2017) &
Ley orgánica. R.O.\ Suplemento No.\ 983, 12-abr-2017 &
Arts.\ 7, 16, 161--168, 175 &
Desarrolla los principios ambientales constitucionales; regula el licenciamiento ambiental, la no regresividad (art.\ 161), el \textit{in dubio pro natura} (art.\ 16) y la responsabilidad objetiva (art.\ 175). \\
\hline

3 &
Reglamento al Código Orgánico del Ambiente (RCOAM) (2019) &
Reglamento. R.O.\ Suplemento No.\ 507, 12-jun-2019. Decreto Ejecutivo No.\ 752 &
Arts.\ 462--463 (declarados inconstitucionales mediante Sentencia 22-18-IN/21) &
Desarrollo reglamentario del COA. Sus arts.\ 462--463 sobre CPLI fueron declarados inconstitucionales por vulnerar la reserva de ley orgánica, conforme a la Sentencia No.\ 22-18-IN/21. \\
\hline

4 &
Ley de Minería (texto consolidado con reformas hasta agosto de 2025) &
Ley ordinaria. R.O.\ Suplemento No.\ 517, 29-ene-2009. Última reforma: R.O.\ Tercer Suplemento No.\ 112, 28-ago-2025 &
Arts.\ 1--4 (objeto y clasificación), 8 (ARCOM), 26 (concesión minera), 55 (contrato de explotación), 78--79 (obligaciones ambientales) &
Eje normativo específico del sector. Define el régimen de concesiones, las categorías de actividad minera según escala y las obligaciones de los concesionarios privados. \\
\hline

5 &
Ley Orgánica Reformatoria a la Ley de Minería (2013) &
Ley orgánica reformatoria. R.O.\ Segundo Suplemento No.\ 37, 16-jul-2013 &
Reformas al régimen de regalías, participación ciudadana y licenciamiento &
Antecedente legislativo. Sus reformas están incorporadas al texto vigente de la Ley de Minería. \\
\hline

6 &
Ley Orgánica para el Fortalecimiento de los Sectores Estratégicos de Minería y Energía (2026) &
Ley orgánica vigente. R.O.\ Quinto Suplemento No.\ 234, 2-mar-2026. Acumula demandas de inconstitucionalidad ante la CCE (al 5-abr-2026: 11 demandas) &
Texto íntegro de la reforma legislativa más reciente &
Objeto central de análisis crítico. Es tratada como norma vigente publicada en R.O.\ Su compatibilidad con el derecho ambiental constitucional es evaluada en el Cap.\ IV. \\
\hline

7 &
Reglamento General a la Ley de Minería (2009) &
Reglamento. R.O.\ Suplemento No.\ 67, 16-nov-2009. Última modificación: 31-ene-2019 &
Procedimientos de concesión, contratos de explotación, obligaciones ambientales reglamentarias &
Desarrollo reglamentario de la Ley de Minería. Complementa el régimen de concesiones y los procedimientos de fiscalización. \\
\hline

8 &
Código Orgánico Integral Penal (COIP) (2014) &
Código. R.O.\ Suplemento No.\ 180, 10-feb-2014 &
Arts.\ 245--254 (delitos contra el ambiente, extracción ilícita de minerales, responsabilidad de personas jurídicas) &
Define los tipos penales de la minería ilegal y los delitos ambientales. Relevante para delimitar el régimen de responsabilidad penal, distinto de la responsabilidad civil objetiva del COA. \\
\hline

9 &
Acuerdo de Escazú (2018) &
Tratado internacional. Ratificado por Ecuador; en vigor desde 22-abr-2021 &
Arts.\ 5--8 (acceso a información ambiental, participación pública, acceso a la justicia) &
Fuente convencional que refuerza las obligaciones de participación y consulta en el sector minero. Complementa el art.\ 398 CRE y los estándares del Convenio OIT 169. \\
\hline

10 &
Instructivo AM No.\ 48 (2015) &
Acuerdo Ministerial. R.O.\ No.\ 637, 27-nov-2015. Última modificación: R.O.\ No.\ 315, 29-ago-2018 &
Procedimientos técnico-administrativos de exploración y explotación minera &
Instrumento de regularización técnica sectorial. Complementa el régimen de licenciamiento específico de la actividad minera privada. \\
\hline

\end{longtable}
\endgroup

\noindent\textit{Nota}: Los instrumentos se presentan en orden de jerarquía normativa y relevancia para la investigación. La columna «Artículos clave» referencia las disposiciones utilizadas en el análisis; no es un listado exhaustivo del contenido normativo de cada instrumento. Los arts.\ 462 y 463 del RCOAM se mencionan únicamente como antecedente del control constitucional por parte de la CCE en la Sentencia No.\ 22-18-IN/21. La Ley Orgánica para el Fortalecimiento de los Sectores Estratégicos de Minería y Energía es tratada como norma vigente, no como proyecto de ley.
